\scanpage{31}
\mspage{30}

% L1–5
\deleted{Thus} Let $\deleted{U}X=\deleted{U}X_{\Phi}$ be a tame subset. An $x\in L$ is said to be in tame relative position with $\deleted{U}X$, if it is so with the $T_i$ (\inserted{\uline{segment pieces}} \uline{constituting} $\deleted{U}X$) i.e. if it is so with $\underset{\partial U}{\underline{\bigcup\partial T_i}}$, or again with $\operatorname{Fr}(\deleted{U}X)$.
% Lleft margin, alongside lines 1–13
% slanted
\margin{More \unsure{generally|generally}, If $\deleted{U}X\unsure{,|\cap}\deleted{\unsure{V|U'}}Y$ are tame, we say that $\deleted{U}X,Y$ are in tame relative position if \deleted{$\partial U$ and $U'$}, $\partial\deleted{U}X,\partial\deleted{V}Y$ are so or again $\operatorname{Fr}\deleted{U}X,\operatorname{Fr}Y$ are so. Example: if $\deleted{U}X\subset Y$, then $(\deleted{U}X,\deleted{V}Y)$ \deleted{\unsure{is}} \unsure{tame|tame}.}
% L5–7
For example \deleted{Let} $x\in\deleted{U}X$, or $x\in\deleted{\unsure{\operatorname{Fr}|\partial}}\partial U$, a fortiori in $\operatorname{Fr}(\deleted{U}X)$, has this property. Let
% L8–9
\[
\begin{aligned}
\deleted{U}X'=C_{\mathrm{mod}}\unsure{U U)|(U)|U)}=\bigl\{x\in L\setminus U\mid{}&x\text{ is in tame relative position}\\
&\text{with respect to }U\bigr\}
\end{aligned}
\]
% L10–13
Then \deleted{\unsure{$\mathcal E$|$\mathcal C$|$C$}} $\deleted{U}X'$ is \unsure{tame|tame}, and it is the largest tame subset $\deleted{V}Y$ of $L$ such that $\deleted{V}Y\subset L\setminus\deleted{U}X$, and $Y$ in tame relative position with respect to $X$. If
% L14
\[
\{T_1\ll T_2<\ldots\ll T_n\}
\]
% L15–20
is the flag defining \deleted{\unsure{$T$|$\Phi$}} $X$, $X'$ is defined by the flag formed of the $n-1$ interstitial segment pieces \deleted{between} for $(T_1,T_2)$, $(T_2,T_3)$ \ldots\ $(T_{n-1},T_n)$, plus (insofar as they exist) the segment pieces, left interstitial of $T_1$, and the right interstitial segment piece of the $n$. \unsure{Continuation|On the other}
